\chapter{插件与技能系统}

\section{插件系统}

Claude Code 的插件系统允许通过第三方扩展来增强功能。

\subsection{架构}

\begin{lstlisting}[language=TypeScript,breaklines=true]
src/plugins/
├── bundled/           # 内置插件
│   └── index.ts       # 内置插件初始化
└── ...

src/services/plugins/
├── pluginLoader.ts    # 插件加载器
└── pluginCliCommands.ts # 插件 CLI 命令

src/utils/plugins/
├── installedPluginsManager.ts # 已安装插件管理
├── managedPlugins.ts          # 托管插件
├── pluginDirectories.ts       # 插件目录
├── cacheUtils.ts              # 缓存工具
├── orphanedPluginFilter.ts    # 孤立插件过滤
└── pluginLoader.ts            # 插件加载逻辑
\end{lstlisting}

\subsection{插件加载}

\begin{lstlisting}[language=TypeScript,breaklines=true]
import { initBuiltinPlugins } from './plugins/bundled/index.js'
import { loadAllPluginsCacheOnly } from './utils/plugins/pluginLoader.js'

// 启动时初始化内置插件
initBuiltinPlugins()

// 缓存优先加载所有插件
const { enabled, errors } = await loadAllPluginsCacheOnly()
\end{lstlisting}

\subsection{插件管理命令}

\begin{lstlisting}[language=TypeScript,breaklines=true]
import {
  VALID_INSTALLABLE_SCOPES,
  VALID_UPDATE_SCOPES,
} from './services/plugins/pluginCliCommands.js'
\end{lstlisting}

支持安装、更新、删除插件，以及按作用域（全局/项目/用户）管理。

\subsection{版本管理}

\begin{lstlisting}[language=TypeScript,breaklines=true]
import { initializeVersionedPlugins } from './utils/plugins/installedPluginsManager.js'
import { cleanupOrphanedPluginVersionsInBackground } from './utils/plugins/cacheUtils.js'
\end{lstlisting}

插件支持版本管理，并在后台清理孤立版本。

\subsection{托管插件}

\begin{lstlisting}[language=TypeScript,breaklines=true]
import { getManagedPluginNames } from './utils/plugins/managedPlugins.js'
\end{lstlisting}

"托管插件"是由 Anthropic 或企业管理员配置的插件，用户无法自行修改。

\section{技能系统}

技能（Skills）是可复用的工作流定义，类似于"预设的操作序列"。

\subsection{架构}

\begin{lstlisting}[language=TypeScript,breaklines=true]
src/skills/
├── bundled/           # 内置技能
│   └── index.ts       # 内置技能初始化
└── ...
\end{lstlisting}

\subsection{内置技能初始化}

\begin{lstlisting}[language=TypeScript,breaklines=true]
import { initBundledSkills } from './skills/bundled/index.js'
\end{lstlisting}

\subsection{技能工具}

\begin{lstlisting}[language=TypeScript,breaklines=true]
import { SkillTool } from './tools/SkillTool/SkillTool.js'
\end{lstlisting}

\code{SkillTool} 是 LLM 调用技能的入口。通过它，模型可以执行预定义的工作流，如 \code{/commit}（提交代码）、\code{/verify}（验证变更）等。

\subsection{技能发现}

\begin{lstlisting}[language=TypeScript,breaklines=true]
// 技能发现预取——每次迭代运行
const pendingSkillPrefetch = skillPrefetch?.startSkillDiscoveryPrefetch(
  null,
  messages,
  toolUseContext,
)
\end{lstlisting}

技能发现在查询循环的每次迭代中运行，与模型流式输出和工具执行并行。这替代了之前在 \code{getAttachmentMessages} 中的阻塞式发现路径。

\subsection{技能搜索}

\begin{lstlisting}[language=TypeScript,breaklines=true]
const clearSkillIndexCache = feature('EXPERIMENTAL_SKILL_SEARCH')
  ? require('./services/skillSearch/localSearch.js').clearSkillIndexCache
  : null
\end{lstlisting}

实验性技能搜索功能支持在大量技能中进行全文搜索。

\subsection{斜杠命令技能}

\begin{lstlisting}[language=TypeScript,breaklines=true]
import { getSlashCommandToolSkills } from './commands.js'
\end{lstlisting}

某些斜杠命令可以注册为技能，使得 LLM 可以自主调用它们。

\section{技能变更检测}

\begin{lstlisting}[language=TypeScript,breaklines=true]
import { skillChangeDetector } from './utils/skills/skillChangeDetector.js'
\end{lstlisting}

当技能文件发生变更时，系统能检测并重新加载，无需重启。

\section{遥测}

\begin{lstlisting}[language=TypeScript,breaklines=true]
import { logSkillsLoaded } from './utils/telemetry/skillLoadedEvent.js'
import { logPluginsEnabledForSession } from './utils/telemetry/pluginTelemetry.js'
import { logPluginLoadErrors } from './utils/telemetry/pluginTelemetry.js'

// 记录加载的技能
void logSkillsLoaded(getCwd(), getContextWindowForModel(model, getSdkBetas()))

// 记录启用的插件
logPluginsEnabledForSession(enabled, managedNames, getPluginSeedDirs())

// 记录插件加载错误
logPluginLoadErrors(errors, managedNames)
\end{lstlisting}

\section{设计启示}

\subsection{技能 vs 工具 vs 命令}

Claude Code 的三层扩展模型：

\begin{longtable}{llll}
\toprule
层次 & 目的 & 调用者 & 粒度 \\
\midrule
\textbf{工具 (Tool)} & 原子操作 & LLM & 低（读文件、执行命令） \\
\textbf{技能 (Skill)} & 工作流 & LLM / 用户 & 中（提交、审查） \\
\textbf{命令 (Command)} & 用户操作 & 用户 & 高（配置、诊断） \\
\bottomrule
\end{longtable}

\subsection{渐进式扩展}

\begin{itemize}[nosep]
  \item 内置插件/技能提供基础功能
  \item 用户可以添加自定义技能
  \item 企业可以通过托管插件统一配置
  \item 第三方可以发布开源插件
\end{itemize}

\bigskip\noindent\rule{\textwidth}{0.4pt}\bigskip

\textbf{下一章}：\href{./chapter-17-patterns.md}{设计模式与工程实践总结}
